
\documentclass[12pt]{article}

\usepackage[utf8]{inputenc}
\usepackage[T1]{fontenc}
\usepackage[spanish]{babel}
\usepackage{amsmath,amssymb}
\usepackage{geometry}
\geometry{margin=2.5cm}

\title{Series de Tiempo y Simulación Estocástica}
\author{Iñaki Castillo}
\date{\today}

\begin{document}

\maketitle

\section{Introducción}

Este documento resume y documenta los conceptos principales desarrollados en los notebooks:
\begin{itemize}
    \item Clase\_1.ipynb
    \item Seasonality.ipynb
    \item Series\_Financieras.ipynb
    \item caminata\_aleat.ipynb
\end{itemize}

\section{Caminata Aleatoria}

Una caminata aleatoria se define como:

\begin{equation}
X_t = X_{t-1} + Z_t
\end{equation}

donde:

\begin{equation}
Z_t \sim N(0,\sigma^2)
\end{equation}

\subsection{Propiedades}

\begin{itemize}
    \item No estacionaria en niveles.
    \item La varianza aumenta con el tiempo.
    \item Es ampliamente utilizada en finanzas.
\end{itemize}

La primera diferencia es:

\begin{equation}
\Delta X_t = X_t - X_{t-1}
\end{equation}

\section{Series Financieras}

Para analizar activos financieros se utilizan retornos logarítmicos:

\begin{equation}
r_t=\ln\left(\frac{P_t}{P_{t-1}}\right)
\end{equation}

donde $P_t$ representa el precio en el periodo actual.

\subsection{Ventajas}

\begin{enumerate}
    \item Son aditivos en el tiempo.
    \item Facilitan el análisis estadístico.
    \item Son estándar en finanzas cuantitativas.
\end{enumerate}

\section{Estacionalidad}

Una serie temporal puede descomponerse como:

\begin{equation}
Y_t = T_t + S_t + E_t
\end{equation}

donde:

\begin{itemize}
    \item $T_t$: Tendencia.
    \item $S_t$: Estacionalidad.
    \item $E_t$: Error aleatorio.
\end{itemize}

\section{Modelo SARIMA}

La notación general es:

\begin{equation}
(p,d,q)(P,D,Q)_m
\end{equation}

donde:

\begin{itemize}
    \item $p$: Autorregresivo.
    \item $d$: Diferenciación.
    \item $q$: Media móvil.
    \item $P$: Autorregresivo estacional.
    \item $D$: Diferenciación estacional.
    \item $Q$: Media móvil estacional.
    \item $m$: Periodicidad.
\end{itemize}

Modelo utilizado:

\begin{equation}
SARIMA(0,1,1)(0,1,1)_{12}
\end{equation}

\section{SARIMAX y Variables Exógenas}

La incorporación de variables externas puede representarse como:

\begin{equation}
Y_t = SARIMA + \beta X_t
\end{equation}

donde $X_t$ representa factores externos relevantes.

\section{Pronósticos}

El objetivo del modelo es obtener:

\begin{equation}
\hat{Y}_{t+h}
\end{equation}

para distintos horizontes de predicción.

\section{Análisis de Residuos}

Los residuos se definen como:

\begin{equation}
e_t = Y_t - \hat{Y}_t
\end{equation}

Un modelo adecuado debe satisfacer:

\begin{equation}
E[e_t] = 0
\end{equation}

y

\begin{equation}
Cov(e_t,e_{t-k}) = 0
\end{equation}

para cualquier rezago $k$.

\section{Conclusiones}

Los notebooks analizados cubren fundamentos de:

\begin{enumerate}
    \item Caminatas aleatorias.
    \item Series financieras.
    \item Estacionalidad.
    \item Modelos SARIMA y SARIMAX.
    \item Pronósticos.
    \item Diagnóstico de residuos.
\end{enumerate}

Estos conceptos constituyen herramientas esenciales para el análisis econométrico y financiero.

\end{document}
